\section{Notation and the tracial completion}\label{r:sec:preliminaries}

For a unital \(C^*\)-algebra \(D\), write \(\T(D)\) for its normalized tracial
states and \(\U(D)\) for its unitary group. For \(\tau\in\T(D)\), put
\[
 \normone{d}{\tau}=\tau(|d|),\qquad
 \normtwo{d}{\tau}=\tau(d^*d)^{1/2}.
\]
These may be seminorms. We use the finite von Neumann algebra obtained
from the tracial GNS representation when taking support projections,
polar decompositions, or signs in a trace calculation. In particular,
tracial Cauchy--Schwarz and trace-norm duality give
\begin{equation}\label{r:eq:trace-inequalities}
 \begin{aligned}
 |\tau(ab)|&\leq\normtwo{a}{\tau}\normtwo{b}{\tau},\qquad
 |\tau(ab)|\leq\norm a\normone b\tau,\\
 \norm{adb}_{s,\tau}&\leq\norm a\norm b\norm d_{s,\tau}
 \qquad(s=1,2).
 \end{aligned}
\end{equation}
In the last expression, \(\norm d_{s,\tau}\) denotes the corresponding
tracial seminorm. All uses of these inequalities also apply to the
unnormalized matrix trace \(\Tr\otimes\tau\). We always reserve \(\Tr\)
for the unnormalized scalar matrix trace. A scalar matrix over a corner
\(qDq\) is written \(G\otimes q\).

A projection is \emph{full} if the closed two-sided ideal it generates is
the whole algebra. For projections \(e,f\in D\), the notation
\(e\precsim f\) means that there is a partial isometry \(v\in D\) with
\(v^*v=e\) and \(vv^*\leq f\). Thus all projection comparisons used below
are implemented inside \(D\).
Real rank zero means that finite-spectrum self-adjoint elements are
norm dense in \(D_{\mathrm{sa}}\).

Every contraction in a unital \(C^*\)-algebra is a linear combination of four
unitaries whose coefficient absolute values sum to at most two. Indeed,
write \(a=h+ik\), where \(h,k\) are self-adjoint contractions, and use
\[
 h=\frac{u_h+u_h^*}{2},\qquad
 u_h=h+i(1-h^2)^{1/2},
\]
and the same formula for \(k\). Consequently, statements about a finite
set of elements can be reduced to statements about a finite set of
unitaries, with controlled constants.

\subsection{The specified completion and limit traces}

Let \(A\) be as in Theorem~\ref{r:thm:main}. For \(a\in A\), write
\[
 \normtwo a{\T(A)}
 =\sup_{\tau\in\T(A)}\tau(a^*a)^{1/2}.
\]
Let \(C\) be the \(C^*\)-algebra of operator-norm-bounded sequences in \(A\)
which are Cauchy for this norm, and let \(J\) consist of those sequences
whose uniform tracial 2-norm tends to zero. The uniform tracial
completion is \(B=C/J\), equipped with its quotient \(C^*\)-norm. The
designated traces on \(B\) are precisely
\[
 \bar\tau([(a_j)])=\lim_j\tau(a_j),\qquad \tau\in\T(A).
\]
These extensions define \(\normtwo{\cdot}{\T(A)}\) on \(B\).
No additional abstract traces of \(B\) enter this notation.

For a fixed free ultrafilter \(\omega\) on \(\N\), define
\begin{equation}\label{r:eq:ultrapower}
 B^\omega=\ell^\infty(\N,B)/I_\omega,\qquad
 I_\omega=\{(b_j):\lim_{j\to\omega}\normtwo{b_j}{\T(A)}=0\}.
\end{equation}
Here and below the sequences in \(I_\omega\) are operator-norm bounded.
For each sequence \((\tau_j)\) in \(\T(A)\), its \emph{limit trace} is
\begin{equation}\label{r:eq:limit-trace}
 \lambda([(b_j)])=\lim_{j\to\omega}\bar\tau_j(b_j).
\end{equation}
Write \(\Lambda_\omega\) for exactly this family of traces.

The copies of \(A\) and \(B\) given by constant sequences are faithful.
For \(A\), each trace is faithful by simplicity. For \(B\), uniform
2-Cauchyness shows that the functions
\(\tau\mapsto\tau(a_j^*a_j)\) converge uniformly. If all their limits
vanish, then \(\normtwo{a_j}{\T(A)}\to0\), so the class in \(B\) is zero.
The same observation identifies the constant copy of \(B\) in
\(B^\omega\).

We say that \(B\) has property \(\Gamma\), or equivalently that \(A\)
has \emph{uniform property \(\Gamma\)}, if there is a projection
\(p\in B^\omega\cap B'\) such that
\begin{equation}\label{r:eq:gamma-definition}
 \lambda(px)=\tfrac12\lambda(x)
 \quad(x\in B,\ \lambda\in\Lambda_\omega).
\end{equation}
This is the tracial-completion formulation of uniform property
\(\Gamma\); see \cite[Definition~16]{STW}. The relation with the standard
all-$m$ formulation uses the following interfaces. The framework of
\cite{CCEGSTW} takes the closed convex hull of \(\Lambda_\omega\) as its
designated ultrapower trace space. For fixed \(p,x\), the identity in
\eqref{r:eq:gamma-definition} extends to this closed convex hull by affine
weak-star continuity, and the supremum tracial 2-norm is unchanged.
The ultrapower identification and transfer
of relative commutants are given in \cite[Proposition~5.7]{CCEGSTW}, and
\cite[Proposition~5.20]{CCEGSTW} identifies uniform property $\Gamma$ of
$A$ with that of its completion. The passage from a halving projection to
equal coefficient-trace partitions is given by
\cite[Proposition~5.23 and Remark~5.24]{CCEGSTW}.
Our proof establishes \eqref{r:eq:gamma-definition} directly for the
specified completion and the originally fixed ultrafilter.
